% 第 6 题：灯泡亮度可调电路（circuitikz）
% 交流电源 U → 开关 S → 触点 a/b/c → 三个串联电阻 R（阶梯，a 在最上端）
% → 原线圈 → 电源，串联回路；副线圈接灯泡 L。匝比 k。
% S 与 a 相连时原线圈回路总阻值为 3R，与 b 为 2R，与 c 为 R。
\begin{tikzpicture}[font=\small, scale=0.9, every node/.style={transform shape=false}]
  % 电源 U（左侧竖直）
  \draw (0,0) to[sV, l=$U$] (0,-2.4);
  % 顶部导线与开关 S
  \draw (0,0) -- (1.5,0);
  \draw (1.5,0) to[switch, l=$S$] (2.9,0);
  \draw (2.9,0) -- (4.4,0);
  % 电阻阶梯（竖直，a 在上端，依次 R、R、R）
  \draw (4.4,0) to[R, l=$R$] (4.4,-1.3);
  \draw (4.4,-1.3) to[R, l=$R$] (4.4,-2.6);
  \draw (4.4,-2.6) to[R, l=$R$] (4.4,-3.9);
  % 触点 a/b/c
  \node[circ, label=left:$a$] at (4.4,0) {};
  \node[circ, label=left:$b$] at (4.4,-1.3) {};
  \node[circ, label=left:$c$] at (4.4,-2.6) {};
  % S 接 a（实线）；接 b、c 的虚线示意
  \draw[dashed] (2.9,0) -- (2.9,-1.3) -- (4.4,-1.3);
  \draw[dashed] (2.9,0) -- (2.9,-2.6) -- (4.4,-2.6);
  % 阶梯下端 → 变压器原线圈（左侧下部）→ 电源下端
  \draw (4.4,-3.9) -- (4.4,-4.6) -- (1.5,-4.6) -- (1.5,-3.6);
  \draw (1.5,-3.6) to[L, l=$k$] (1.5,-2.4);
  \draw (1.5,-2.4) -- (0,-2.4);
  % 副线圈与灯泡 L（右侧）
  \draw (2.0,-3.6) to[L] (2.0,-2.4);
  \draw (2.0,-3.6) -- (3.6,-3.6) to[lamp, l=$L$] (3.6,-2.4) -- (2.0,-2.4);
  % 铁芯（两线圈之间）
  \draw[thick] (1.72,-3.7) -- (1.72,-2.3);
  \draw[thick] (1.78,-3.7) -- (1.78,-2.3);
\end{tikzpicture}
